\lesson{Extensions to the object calculus}\label{sec:larger}
\focus{Fields, recursive types, and intersections introduce additional evidence and normalization cases.}

\subsection*{Field presence and field types}
An object can hold a field $a$ as well as define a type $A$.
A bound on the field's type does not by itself prove that the field
exists. The target also associates type names with field labels:
$s.a$ names the type of field $a$'s body. This is a \emph{block name},
like $s.A$ for a type member; here it denotes a type, not the term that
projects the field's value.
The target therefore records two separate propositions:
a presence fact for $a$, and an inclusion from $s.a$ to its advertised
type.

A literal's precise information includes the field's actual type
definition and its presence. Presence evidence lets the safety proof
establish that field lookup succeeds. Projection runs the stored field
body, and inclusion evidence presents its result at the interface's
declared type.
As with type members, the interface records propositions justified by
the stored literal.

For example, consider an object defining $A=W$ and holding
$a=f$. A client may see its field at type $x.A$ even if the stored field
was originally checked at $W$. The precise facts $s.a\equiv W$ and
$s.A\equiv W$ connect the field's type with the type member.
At $x$, these facts justify presenting the field result at $x.A$.
Field projection and the coercion of its result are separate operations,
each with explicit evidence.

\subsection*{Recursive self types}
A source self type can mention its own members:
an object might advertise that its field $a$ has type $s.A$.
Opening this interface at $x$ replaces the bound self by $x$.
Closing it again restores the self binder.
The target uses fold and unfold wrappers on atoms to record these
changes of presentation. Like casts, the wrappers preserve the root.

The self binder in a telescope allows propositions to refer to their
object. Recursive source typing additionally requires rules for opening
and closing the interface at a particular root.

\subsection*{Intersections and self-bounds}
An intersection $S\wedge T$ describes a value usable at both $S$ and $T$.
The field interface $\{a:T\}$ promises a field $a$ with type $T$.
When both components list member declarations, their intersection can
be represented by collecting facts from both lists.
But the source also allows $x.A\wedge\{a:T\}$, where the selection
$x.A$ supplies no visible list of declarations.

To handle components without a visible telescope, the translation
introduces a \emph{self-bound}: a proposition
saying that the whole self has some type. In this example, the $x.A$
component contributes the fact that the whole self can be used at $x.A$;
the field component contributes presence and a type bound for $a$.

Self-bounds require additional head forms and entries that record
conversions through a bound to the interface used by an operation.
These cases extend the five forms in Section~\ref{sec:normalization};
their typing still ensures that normalized evidence supplies the
information required by the operation.

\subsection*{Cyclic aliases}
An object's type definitions may say $A=s.B$ and $B=s.A$.
Unbounded traversal of these definitions would diverge. The resolver
uses a finite bound derived from the available definitions and resolves
a cycle to $\top$, which contributes no additional information.
Resolving it to $\bot$ would instead assert an uninhabited value type.
This result describes the resolved head shape; it supplies neither a
source typing rule nor equality evidence between a cyclic name and
$\top$.

The bounds for alias resolution and evidence normalization have distinct
justifications. Alias resolution uses finiteness of the available names;
evidence normalization uses induction on typing to establish the
existence of sufficient fuel.

\subsection*{Related developments}
The repository also contains a paths variant, where a field path such as
$x.a.b$ can be the receiver of a type selection $(x.a.b).A$.
Singleton types there record that a value is the one named by a
particular path.
A separate OOPSLA 2016 development has recursive subtyping under a self
assumption and its own evidence target.
The theorems apply to these developments separately; they do not cover
their unrestricted combination. Such a combination would need to
establish field presence, stable path resolution, and recursive
subtyping evidence under the appropriate self assumptions within one
evidence language and machine semantics.
